% TOP_PROBLEM: {"id": 2758, "title": "Kirby Problem 2.10", "category_id": 7, "category": {"id": 7, "name": "topology", "display_name": "Topology", "description": "Properties preserved under continuous deformations.", "slug": "topology", "order_index": 7, "created_at": "2026-07-31T15:26:25.671Z"}, "status": "open", "problem_number": "KP-2.10", "set_id": 11}
% FIRST_POSTED: 2026-10-01
% ENTRY_KIND: statement_only
% UnsolvedMath ID 2758; source catalogue code KP-2.10
% !TeX program = lualatex
% Definition and sources only; this file is not an LLM solution attempt.
\documentclass[11pt,a4paper]{article}
\usepackage[margin=25mm]{geometry}
\usepackage{fontspec}
\setmainfont{lmroman10-regular.otf}[BoldFont=lmroman10-bold.otf,
 ItalicFont=lmroman10-italic.otf,BoldItalicFont=lmroman10-bolditalic.otf]
\usepackage{amsmath,amssymb,mathtools}
\usepackage{unicode-math}
\setmathfont{latinmodern-math.otf}
\AtBeginDocument{\let\setminus\smallsetminus}
\usepackage{xurl,hyperref}
\hypersetup{colorlinks=true,urlcolor=blue}
\setcounter{secnumdepth}{0}
\setlength{\parindent}{0pt}
\setlength{\parskip}{5pt}
\setlength{\emergencystretch}{3em}
\begin{document}
\section*{Kirby Problem 2.10}\label{2758}
{\small UnsolvedMath ID 2758 · KP-2.10 · Topology · Kirby's Problems in Low-Dimensional Topology}
\subsection{Definitions and mathematical statement}
Let $S$ be a finite-type oriented surface of complexity $3g-3+p\geq1$, where
$g$ is its genus and $p$ its number of punctures. Its curve graph $\mathcal C(S)$
has one vertex for each isotopy class of essential nonperipheral simple closed
curve, and, at complexity at least two, an edge for disjoint representatives. Every edge
has length one; $d_{\mathcal C}$ is shortest-path distance. The customary
minimal-intersection adjacency is used for the once-punctured torus and
four-punctured sphere.

(a) Find an efficient algorithm computing the \emph{exact} distance between two
input curves. For a uniform complexity question, encode $S$ by a triangulation
and curves by normal coordinates in binary; running time should be controlled
in the total input length, including surface complexity, rather than only for
one fixed surface.
Ask the analogous question for the arc graph on a surface with boundary:
vertices are isotopy classes of essential properly embedded arcs, with
endpoints allowed to move along boundary components, and edges mean
representatives with disjoint interiors.

(b) Find algorithms computing distances between effectively described
quasiconvex subsets $A,B\subseteq\mathcal C(S)$, where
$d(A,B)=\min\{d(a,b):a\in A,b\in B\}$.
A subset is quasiconvex if geodesics between its vertices stay within some
bounded distance of it. A principal case is the disk set of a handlebody:
boundary curves bounding properly embedded disks. For a Heegaard splitting
$H_1\cup_S H_2$, compute the distance between its two disk sets from finite
handlebody and gluing data. These infinite sets require specified effective
descriptions; arbitrary unencoded subsets are not algorithmic inputs.
\subsection{Short English statement}
Compute exact distances in curve graphs efficiently, and develop algorithms for distances between naturally specified infinite subsets, especially the disk sets of Heegaard splittings.
\subsection{Sources}
\begin{itemize}
\item[S1] ULAM.AI and the UnsolvedMath Contributors, supplied problems.json, record 2758 (KP-2.10), Kirby's Problems in Low-Dimensional Topology. Catalogue identity and source transcription; CC BY 4.0.
\url{https://huggingface.co/datasets/ulamai/UnsolvedMath}.
\item[S2] K3: A New Problem List in Low-Dimensional Topology, Problem 2.10. Expanded formulation of the supplied catalogue statement.
\url{https://math.berkeley.edu/sites/default/files/surv-295-ruberman-watermarked-author-pdf.pdf}.
\end{itemize}
\end{document}
